% Section 3 — Model description.
% Largely ported from the ECCV Method section, with terminology
% adjusted for the GMD audience and equation references kept stable.

This section summarises the CRAN-PM model architecture. The
mathematical motivation and detailed comparison to alternative
designs are reported in the methodology paper~\citep{kheder2026eccv};
here we focus on the level of detail required to read the source
code and reproduce a forecast.

\subsection{Forecasting problem and conventions}
\label{sec:model:problem}

We forecast surface PM\textsubscript{2.5} at 1\,km spatial resolution
over Europe ($30^\circ$N--$72^\circ$N, $25^\circ$W--$45^\circ$E,
$4{,}192 \times 6{,}992$ pixels at 0.01$^\circ$) at lead times
$\tau \in \{1, 2, 3, 4\}$ days. The forecast valid at day $t+\tau$
is decomposed as
\begin{equation}
\hat{y}_{t+\tau}(s) = x^{\ell}_t(s) + f_\theta\!\big(\mathbf{x}^{\text{g}}_t,
\mathbf{x}^{\ell}_t, \tau\big)(s),
\label{eq:delta}
\end{equation}
where $x^{\ell}_t(s)$ is the observed (GHAP) PM\textsubscript{2.5}
at pixel $s$ on day $t$, $\mathbf{x}^{\text{g}}_t$ is the global
meteorology + chemistry input and $\mathbf{x}^{\ell}_t$ the local
high-resolution input. The neural network $f_\theta$ outputs a
\emph{delta} (residual change) that is added to today's analysis.
The final convolution of $f_\theta$ is initialised to zero, so that
training starts from a persistence baseline $\hat{y}_{t+\tau} \equiv
x^{\ell}_t$ and the model only needs to learn the day-to-day
\emph{change} \citep{bi2023pangu, lam2023graphcast}.

\subsection{Inputs}
\label{sec:model:inputs}

\paragraph{Global branch (25\,km).}
The global input $\mathbf{x}^{\text{g}}_t \in \mathbb{R}^{70 \times 168 \times 280}$
stacks (i)~ERA5 surface and pressure-level fields at days $t$ and
$t-1$ (60 channels), and (ii)~CAMS analysis of six pollutants
(NO\textsubscript{2}, O\textsubscript{3}, SO\textsubscript{2}, CO,
PM\textsubscript{10} plus PM\textsubscript{2.5}) at days $t$ and
$t-1$ (10 channels). The grid is 0.25$^\circ$. Variable definitions,
units and pressure levels are listed in Table~\ref{tab:inputs}.

\paragraph{Local branch (1\,km).}
The local input $\mathbf{x}^{\ell}_t \in \mathbb{R}^{5 \times 512 \times 512}$
contains, on a $512\!\times\!512$ tile, (i)~the GHAP
PM\textsubscript{2.5} analysis at days $t$ and $t-1$,
(ii)~GMTED2010 elevation downsampled to 0.01$^\circ$,
(iii)~latitude and longitude as auxiliary positional channels.
A pan-European inference pass tiles the domain into 126 overlapping
tiles and stitches the residual outputs.

\subsection{Architecture overview}
\label{sec:model:arch}

Figure~\ref{fig:architecture} shows the dual-branch design.
\begin{enumerate}
\item The \textbf{global branch} encodes coarse meteorology with a
Vision Transformer of patch size $8 \times 8$, embedding dimension
$d_g = 768$, depth 8 (one elevation-aware block followed by 7
Swin-style blocks) and 12 attention heads. Patches are reordered
along the local wind direction prior to the transformer
(\emph{wind-guided shuffling}, \S\ref{sec:model:wind}).

\item The \textbf{local branch} processes a $512 \times 512$ tile
through a Vision Transformer of patch size $16 \times 16$,
$d_\ell = 512$, depth 6 and 8 heads. A learnable lead-time
embedding $\mathbf{E}_{\text{lead}}(\tau)$ is added to the patch
embeddings.

\item A \textbf{cross-resolution attention bridge}
(\S\ref{sec:model:cross}) lets local tokens query the global tokens
and integrate large-scale context.

\item A \textbf{convolutional decoder} progressively upsamples the
fused $32 \times 32$ token grid to $512 \times 512$ via four
PixelShuffle blocks \citep{shi2016pixelshuffle}, ending with a
zero-initialised $1 \times 1$ convolution that produces the residual
$\Delta$ (\S\ref{sec:model:problem}).
\end{enumerate}

\subsection{Elevation-aware attention}
\label{sec:model:elev}

For an attention pair $(i, j)$ with mean tile elevations $E_i, E_j$,
we add an asymmetric bias to the attention logits before the
softmax:
\begin{equation}
B^{\text{elev}}_{ij} = -\alpha \,\text{ReLU}\!\Big(\frac{E_j - E_i}{E_0}\Big),
\quad B^{\text{elev}}_{ij} \in [-10, 0],
\end{equation}
with $\alpha$ learnable per head and $E_0 = 1{,}000$\,m for the
global branch and $500$\,m for the local branch. The bias penalises
attention from low-elevation queries to high-elevation keys, which
is consistent with the dominant katabatic transport pattern in
mountain valleys \citep{whiteman2000mountain}.

\subsection{Wind-guided cross-attention}
\label{sec:model:cross}

The cross-resolution bridge is composed of two cross-attention
layers in which local tokens act as queries and the global tokens,
linearly projected to dimension $d_\ell$, act as keys and values.
A wind-guided bias encodes upwind alignment:
\begin{equation}
\gamma_{ij} = \frac{(\mathbf{p}_i - \mathbf{p}_j)}{\|\mathbf{p}_i -
\mathbf{p}_j\|} \cdot \frac{(u_{10,i}, v_{10,i})}{\|(u_{10,i}, v_{10,i})\|},
\qquad
B^{\text{wind}}_{ij} = \beta \cdot \gamma_{ij},
\end{equation}
where $\mathbf{p}_i$ are tile-centre coordinates, $(u_{10}, v_{10})$
is the surface wind from ERA5 and $\beta$ is learnable per head.
Adding $B^{\text{wind}}$ to the cross-attention logits makes the
bridge prefer to attend in the upwind direction, which is the
physical advection path of pollutants.

\subsection{Wind-guided shuffling in the global branch}
\label{sec:model:wind}

Inside the global branch, each $7 \times 7$ patch group is
reordered before the transformer according to a 16-sector
quantisation of the local mean wind. This aligns the sequence
processing with the transport direction \citep{topoflow}.

\subsection{Training objective}
\label{sec:model:loss}

The model is trained with a sum of three terms:
\begin{equation}
\mathcal{L} = \mathcal{L}_{\text{pixel}}
            + \lambda_{\text{FFL}} \,\mathcal{L}_{\text{FFL}}
            + \lambda_{\text{station}} \,\mathcal{L}_{\text{station}}.
\end{equation}
$\mathcal{L}_{\text{pixel}}$ is a land-only mean-squared error
between $\hat{y}_{t+\tau}$ and the GHAP target.
$\mathcal{L}_{\text{FFL}}$ is the focal frequency loss
\citep{jiang2021focal} that emphasises high-frequency spatial
errors. $\mathcal{L}_{\text{station}}$ anchors the prediction at
the $N_s = 2{,}971$ EEA stations included in the training set,
acting as a soft regulariser against systematic biases.
We set $\lambda_{\text{FFL}} = \lambda_{\text{station}} = 0.1$.
